\section{The \texttt{Incoherent} McStas Component}
Incoherent sample (such as Vanadium) sample, with quasielastic component OR or global energy transfer.

\subsection*{Identification}
\begin{itemize}
  \item \textbf{Author:} Kim Lefmann and Kristian Nielsen
  \item \textbf{Origin:} Risoe
  \item \textbf{Date:} 15.4.98
\end{itemize}

\subsection*{Description}
A Double-cylinder shaped incoherent scatterer (like Vanadium) with both elastic and quasielastic (Lorentzian) components. No multiple scattering (but approximation available). Absorption included. \textbf{Sample focusing:} The area to scatter to is a disk of radius 'focus\_r' situated at the target. This target area may also be rectangular if specified focus\_xw and focus\_yh or focus\_aw and focus\_ah, respectively in meters and degrees. The target itself is either situated according to given coordinates (x,y,z), or defined with the relative target\_index of the component to focus to (next is +1). This target position will be set to its AT position. When targeting to centered components, such as spheres or cylinders, define an Arm component where to focus to. \textbf{Sample shape:} Sample shape may be a cylinder, a sphere, a box or any other shape

\begin{verbatim}
box/plate:       xwidth x yheight x zdepth (thickness=0)
\end{verbatim}

hollow box/plate:xwidth x yheight x zdepth and thickness\textgreater{}0

\begin{verbatim}
cylinder:        radius x yheight (thickness=0)
\end{verbatim}

hollow cylinder: radius x yheight and thickness\textgreater{}0

\begin{verbatim}
sphere:          radius (yheight=0 thickness=0)
hollow sphere:   radius and thickness>0 (yheight=0)
any shape:       geometry=OFF file
\end{verbatim}

The complex geometry option handles any closed non-convex polyhedra. It computes the intersection points of the neutron ray with the object transparently, so that it can be used like a regular sample object. It supports the PLY, OFF and NOFF file format but not COFF (colored faces). Such files may be generated from XYZ data using: qhull \textless{} coordinates.xyz Qx Qv Tv o \textgreater{} geomview.off or powercrust coordinates.xyz and viewed with geomview or java -jar jroff.jar (see below). The default size of the object depends of the OFF file data, but its bounding box may be resized using xwidth,yheight and zdepth.

Example: Incoherent(radius=0.05,focus\_r=0.035, pack=1, target\_index=1) Incoherent(geometry="socket.off",focus\_r=0.035, pack=1, target\_index=1)

\subsection*{Input parameters}
Parameters in \textbf{boldface} are required; the others are optional.

\begin{longtable}{p{0.22\textwidth}p{0.12\textwidth}p{0.46\textwidth}p{0.14\textwidth}}
\toprule
\textbf{Name} & \textbf{Unit} & \textbf{Description} & \textbf{Default} \\
\midrule
\endhead
geometry & str & Name of an Object File Format (OFF) or PLY file for complex geometry. The OFF/PLY file may be generated from XYZ coordinates using qhull/powercrust & 0 \\
radius & m & Outer radius of sample in (x,z) plane & 0 \\
xwidth & m & Horiz. dimension of sample (bounding box if off file), as a width & 0 \\
yheight & m & Vert.  dimension of sample (bounding box if off file), as a height. A sphere shape is used when 0 and radius is set & 0 \\
zdepth & m & Depth of sample (bounding box if off file) & 0 \\
thickness & m & Thickness of hollow sample & 0 \\
target\_x &  &  & 0 \\
target\_y & m & position of target to focus at & 0 \\
target\_z &  &  & 0 \\
focus\_r & m & Radius of disk containing target. Use 0 for full space & 0 \\
focus\_xw & m & horiz. dimension of a rectangular area & 0 \\
focus\_yh & m & vert.  dimension of a rectangular area & 0 \\
focus\_aw & deg & horiz. angular dimension of a rectangular area & 0 \\
focus\_ah & deg & vert.  angular dimension of a rectangular area & 0 \\
target\_index & 1 & Relative index of component to focus at, e.g. next is +1 & 0 \\
pack & 1 & Packing factor & 1 \\
p\_interact & 1 & MC Probability for scattering the ray; otherwise transmit & 1 \\
f\_QE & 1 & Fraction of quasielastic scattering (rest is elastic) & 0 \\
gamma & 1 & Lorentzian width of quasielastic broadening (HWHM) & 0 \\
Etrans & meV & Global energy-transfer, for use in inelastic settings & 0 \\
deltaE & meV & Width in energy around Etrans, for use in inelastic settings & 0 \\
sigma\_abs & barns & Absorption cross section pr. unit cell at 2200 m/s & 5.08 \\
sigma\_inc & barns & Incoherent scattering cross section pr. unit cell & 5.08 \\
Vc & \AA{}$^{3}$ & Unit cell volume & 13.827 \\
concentric & 1 & Indicate that this component has a hollow geometry and may contain other components. It should then be duplicated after the inside part (only for box, cylinder, sphere) & 0 \\
order & 1 & Limit multiple scattering up to given order 0 means all (default), 1 means single, 2 means double, ... & 0 \\
\bottomrule
\end{longtable}

\subsection*{Links}
\begin{itemize}
  \item Component source code found in file \texttt{Incoherent.comp}.
  \item \htmladdnormallink{Cross sections for single elements}{http://www.ncnr.nist.gov/resources/n-lengths/}
  \item \htmladdnormallink{Cross sections for compounds}{}
  \item \htmladdnormallink{Web Elements}{http://www.webelements.com/}
  \item \textless{}A HREF="http://neutron.risoe.dk/mcstas/components/tests/Incoherent/"\textgreater{}Test
  \item results\textless{}/A\textgreater{} (not up-to-date).
  \item The test/example instrument \htmladdnormallink{vanadium\_example.instr}{../examples/vanadium\_example.instr}.
  \item The test/example instrument \htmladdnormallink{QENS\_test.instr}{../examples/QENS\_test.instr}.
  \item \htmladdnormallink{Geomview and Object File Format (OFF)}{http://www.geomview.org}
  \item Java version of Geomview (display only) \htmladdnormallink{jroff.jar}{http://www.holmes3d.net/graphics/roffview/}
  \item \htmladdnormallink{qhull}{http://qhull.org}
  \item \htmladdnormallink{powercrust}{http://www.cs.ucdavis.edu/\textasciitilde{}amenta/powercrust.html}
\end{itemize}
\IfFileExists{samples/Incoherent_static.tex}{\input{samples/Incoherent_static.tex}}{}